\documentclass[11pt]{article}

\usepackage[margin=1in]{geometry}
\usepackage{amsmath, amssymb, amsfonts}
\usepackage{setspace}
\usepackage{hyperref}

\setstretch{1.15}

\title{Methods Appendix:\\Estimating Local LGBTQ Populations}
\author{ }
\date{\today}

\begin{document}
\maketitle

\section*{Overview}

I estimate the expected LGBTQ population of every U.S. Census tract in five
conceptual stages: (1) measure LGBTQ identification and gender-identity
composition in Household Pulse Survey (HPS) microdata, softly calibrated to
external Gallup benchmarks; (2) stabilize these rates against sampling noise using two complementary smoothing techniques — one across nearby ages, one across states; (3) project the resulting rates onto Census tracts using ACS 5-year age–sex population counts; (4) reweight the within-state spatial distribution using the observed geographic concentration of same-sex couples and singles within an age-sex cell as a proxy for the broader LGBTQ population, smoothed across space so the result doesn't show artificial jumps at PUMA boundaries; and (5) enforce a capacity constraint so no tract's estimate can exceed its underlying population. Each stage is detailed in its own section below, with the precise notation defined immediately below.

\section*{Notation}

Throughout, $i$ indexes survey respondents and $w_i$ denotes HPS survey
weights. I fix notation of gender and sex up front:

\begin{itemize}
  \item \textbf{Sex assigned at birth}, $s \in \{M, F\}$. Since the ACS does not
  ask about gender identity, I begin with sex assigned at birth.
  Table B01001 provides tract-level counts of age-sex cells.
  \item \textbf{Gender identity}, $g \in \{\text{cis}, \text{nb},
  \text{trans}\}$, estimated from HPS in Section~\ref{sec:dirichlet}. This
  is \emph{not} the same as the final output gender bucket used for
  population accounting on the map. Instead, the gender identity variable
  simply flags an individual's relationship to their sex assigned at birth.
  For example, transgender respondent's gender identity is not their sex
  assigned at birth, so ``trans'' respondents with
  $s=M$ are trans women and those with $s=F$ are trans men.
\end{itemize}

I therefore define a third symbol, the \textbf{gender bucket}
$k \in \{m, w, \text{nb}\}$ (man, woman, non-binary), combining the sex
assigned at birth and gender identity:
\[
k(s,g) =
\begin{cases}
  m  & (s=M,\ g=\text{cis}) \ \text{or}\ (s=F,\ g=\text{trans}) \\
  w  & (s=F,\ g=\text{cis}) \ \text{or}\ (s=M,\ g=\text{trans}) \\
  \text{nb} & g = \text{non-binary}
\end{cases}
\]
In words: a cisgender respondent's bucket matches their sex assigned at
birth; a transgender respondent's bucket does not match their
sex assigned at birth.  Non-binary respondents form their own bucket
regardless of $s$. Sections~\ref{sec:acs}--\ref{sec:capacity} report
quantities indexed by $k$; Section~\ref{sec:dirichlet} reports the
intermediate quantities indexed by $g$ that $k$ is built from.

The remaining recurring symbols are: $a$ for exact age or an ACS age bin;
$\ell \in \{\text{LG}, \text{Bisexual},
\text{Queer}\}$ for LGBT subcategory (excluding Trans, which is cross-cut
with $k$ rather than treated as a fourth member of $\ell$); $p$ for a PUMA
(Public Use Microdata Area); and $b$ for a Census tract.

\section{Measurement in HPS Microdata}

Let $Y_i \in \{0,1\}$ indicate whether individual $i$ identifies as LGBTQ.
In the HPS, I mark respondents as LGBTQ if they report a sexual orientation besides straight, a trans
gender identity, or both. Note that non-binary people are not automatically LGBTQ,
though an overwhelming majority report a non-straight sexuality too.
Sex assigned at birth $s_i$ and age $a_i$ are taken directly from HPS
(age top-coded at 62 to limit disclosure risk in the sparse cells of older cohorts).
Survey weights $w_i$ are retained throughout all of the estimation steps
below.

\section{Soft Calibration to Gallup Marginals}

HPS and Gallup differ in survey mode, sampling variation, and field timing,
so I don't expect them to agree exactly and I don't force them to.
Instead, HPS-based probabilities are \emph{softly} calibrated toward Gallup
national benchmarks, nudged toward external validity.

Each individual begins with a baseline probability $p_{0i} =
\text{clipped}(Y_i)$ (clipped away from exactly 0 or 1 so the subsequent
logit transform is well-defined). A logit adjustment model is estimated:
\[
\text{logit}(p_i) = \text{logit}(p_{0i}) + X_i \theta,
\]
where $X_i$ includes generation effects and generation-by-sex interactions.
The parameter vector $\theta$ minimizes a weighted objective combining
(i) a deviation penalty that keeps estimates close to HPS baselines, and
(ii) a margin penalty that discourages deviation from Gallup targets,
optimized via BFGS. In practice, this keeps the vast majority of
state-level estimates within one to two percentage points of Gallup's
national top-line estimate of roughly 9 percent, while still letting
genuine state-level heterogeneity show through.

\section{State--Sex--Age LGBT Rate Estimation}

Using calibrated probabilities, sex--age--state LGBT identification rates
are computed as a straightforward weighted mean:
\[
\widehat{P}(Y=1 \mid s, a, \text{state})
=
\frac{\sum_i w_i\, p_i^{\text{adj}}}{\sum_i w_i},
\]
with effective sample size $n_{\text{eff}} = (\sum_i w_i)^2 / \sum_i
w_i^2$ tracking how much genuine information backs each cell. Weekly HPS
survey waves are pooled and averaged to reduce short-run noise before this
step.

\section{Age-Local Dirichlet Smoothing of Gender Composition}
\label{sec:dirichlet}

This step estimates $P(g \mid \ell, s, a)$ -- the distribution of
gender-identity type $g \in \{\text{cis}, \text{nb}, \text{trans}\}$ within
each LGBT subcategory, sex assigned at birth, and exact age (see
\emph{Notation} above for why $g$ is deliberately kept distinct from the
final output bucket $k$).

The raw cell rate is
\[
\hat{p}_{g,\ell,s,a} = \frac{n_{g,\ell,s,a}}{\sum_g n_{g,\ell,s,a}},
\qquad
n_{\text{eff},\ell,s,a} = \frac{(\sum w_i)^2}{\sum w_i^2},
\]
where $n_{g,\ell,s,a}$ denotes survey-weighted counts. Some cells are thin;
in particular, non-binary and transgender respondents at a single exact age within
a single sexuality category might often yield sparse cells. 
I address this with two smoothing passes: first
borrowing strength from \emph{nearby ages} within a state (this section),
then, in Section~\ref{sec:crossstate}, pooling the coarser LGBT subcategory
shares \emph{nationally}.

\subsection{Age-Local Prior Construction}

I motivate the smoothing with a concrete pattern in the data: younger
AFAB (assigned female at birth) lesbians are roughly 15 percentage points
more likely to identify as non-binary than AMAB gay men in the same age
cohort. I also don't want a single noisy exact-age cell to dominate the
estimate. Age-local smoothing does exactly this: it borrows
information from \emph{nearby} ages, weighted by how much data backs each
neighbor and how close in age they are.

For a focal age $a$, neighboring ages $a' \neq a$ receive weights
\[
w_{a'}^{(a)} =
n_{\text{eff},\ell,s,a'}\,
\exp\!\left(-\frac{|a-a'|}{\tau}\right),
\]
where $\tau > 0$ controls how quickly influence decays with age distance.
The age-local prior mean is then the weighted average of neighboring rates:
\[
m_g^{(a)} =
\frac{\sum_{a' \neq a} w_{a'}^{(a)}\, \hat{p}_{g,\ell,s,a'}}
{\sum_{a' \neq a} w_{a'}^{(a)}}.
\]

\subsection{Dirichlet Moment Matching}

A Dirichlet prior is centered at $\mathbf{m}^{(a)}$, with concentration
parameter estimated via method of moments:
\[
\alpha_0 =
\text{mean}_g\!\left(\frac{m_g(1-m_g)}{v_g} - 1\right),
\]
where $v_g$ is the weighted variance of $\hat p_{g,\ell,s,a'}$ across
neighboring ages $a'$ -- intuitively, how much the neighbors agree with
each other. The prior vector is $\alpha_g = \alpha_0\, m_g$.

\subsection{Posterior Shrinkage}

The smoothed estimate for focal age $a$ blends the raw cell rate with the
age-local prior, weighted by how much real data each side has:
\[
\tilde{p}_{g,\ell,s,a}
=
\frac{n_{\text{eff},\ell,s,a}\,\hat{p}_{g,\ell,s,a} + \alpha_g}
{n_{\text{eff},\ell,s,a} + \alpha_0}.
\]
A well-sampled cell ($n_{\text{eff}}$ large relative to $\alpha_0$) stays
close to its raw rate; a thin cell shrinks toward the age-local prior.

\section{National LGBT Subcategory Composition}
\label{sec:crossstate}

For each sex $s$ and age bin $a$, I also need the split between LG, Bisexual, and Queer identification after conditioning on LGBTQ identity. I pool nationally instead of estimating by state:
\[
\tilde\pi_{\ell,s,a}
=
\frac{\hat\pi_{\ell,s,a}}{\sum_{\ell'} \hat\pi_{\ell',s,a}},
\qquad
\hat\pi_{\ell,s,a} = \text{mean over HPS weeks of the national }
(s,a)\text{ share identifying as } \ell.
\]
Every state now applies the same LG/Bisexual/Queer mix to its own
state-level, Gallup-calibrated LGBT identification rate $\widehat
P(Y=1\mid s,a,\text{state})$ from Section 3 -- states differ in \emph{how
many} people identify as LGBT, but not in \emph{the mix
of labels} people use.

\section{ACS Projection to Census Tracts}
\label{sec:acs}

Let $N_b^{(s,a)}$ denote the ACS 5-year population count for sex assigned
at birth $s \in \{M,F\}$ and age bin $a$ in Census tract $b$. I now combine
every rate estimated above -- identification, subcategory share, and
gender-identity composition -- with local ACS demographics to get an
expected count for each tract, LGBT subcategory, and output gender bucket
$k \in \{m, w, \text{nb}\}$:
\[
\hat{E}[\text{count}_{k,\ell,b}]
=
\sum_{s,a} N_b^{(s,a)}
\cdot \widehat{P}(Y=1 \mid s,a,\text{state})
\cdot \tilde\pi_{\ell,s,a}
\cdot \!\!\sum_{g\,:\,k(s,g)=k}\!\! \tilde p_{g,\ell,s,a}.
\]
The innermost sum collapses the gender-identity-type rate $\tilde
p_{g,\ell,s,a}$ from Section~\ref{sec:dirichlet} into the output bucket $k$
using the $(s,g)\mapsto k$ mapping defined in \emph{Notation}. Because
identification, subcategory, and gender-identity rates all vary with age
and sex, and both vary locally, expected LGBTQ counts vary across tracts
even within a single state -- this section captures only the
\emph{demographic composition} channel of that variation; the next section
adds a \emph{geographic clustering} channel on top of it. For the PUMA
reweighting step below, the per-cell summand of this sum -- indexed by
$(s,g,\ell,a,b)$, before summing over $s,a$ and collapsing $g$ into $k$ --
is what actually gets reweighted; this formula shows the result after that
collapse.

Tracts that span PUMA boundaries receive population-weighted contributions
from each overlapping PUMA, proportional to the area of intersection.

\section{PUMA Reweighting}

Spatial sorting shapes the spatial distribution of LGBTQ people across the
United States. LGBTQ people cluster for many reasons, including proximity
to community infrastructure, local social tolerance, idiosyncratic amenity
preferences, and more. Yet the projection in Section~\ref{sec:acs}
allocates LGBTQ population using only local age and sex composition -- it
offers no mechanism for the spatial sort of LGBTQ people beyond age and
sex. I use the observed concentration of same-sex couples as a proxy for
the broader LGBTQ population's spatial distribution.

Same-sex couples are an imperfect proxy, skewing older,
wealthier, better-educated, and more likely to identify as gay or lesbian
than the broader LGBTQ population. Within metropolitan areas, their spatial 
Concentration might systematically differ other LGBTQ sub-groups, particularly young, non-binary and transgender populations. Rather than
apply a couples-based tilt to the whole LGBTQ population, I blend it with a second spatial proxy: the geographic concentration of
\emph{singles} of the same sex-age cell. The weight placed on couples relative to single depends on the national age-sex cell marriage rate in the HPS.

\subsection{Coupled and Singles Spatial Signals}

I construct a same-sex-couple signal from ACS Table B09019 counts of
same-sex (SS) and opposite-sex (OS) couples by PUMA, unchanged from before:
PUMAs with fewer than 100 total couples are dropped as too noisy to trust,
and the raw same-sex share is winsorized to the 1st--99th percentile across
PUMAs before transforming, to keep a handful of tiny, extreme PUMAs from
dominating the calibration,
\[
\hat q_p = \text{winsorize}_{[0.01,\,0.99]}
\!\left(\frac{\text{SS couples}_p}{\text{SS couples}_p + \text{OS
couples}_p}\right).
\]
I normalized to a spatial share:
\[
\rho_p^{\text{couple}} = \frac{\text{SS couples}_p}{\sum_{p'} \text{SS
couples}_{p'}}.
\]

The singles signal comes from ACS Table B12002 (sex by marital status
by age), for birth sex $s$ and age bin $a$:
\[
\rho_p^{\text{single}}(s,a) = \frac{\text{Not-married}_{p,s,a}}
{\sum_{p'} \text{Not-married}_{p',s,a}},
\]
where ``not married'' includes everyone except people currently married with a
spouse present. PUMAs with fewer than 100 such people in a
given $(s,a)$ are dropped, the same threshold applied to the couples signal.

\subsection{Coupled-Share Blend Weight}

From the HPS, I estimate the national married share of people in cell $(s,g,\ell,a)$.
HPS lacks an option for unmarried cohabiting partners question, so ``married'' is an imperfect proxy for ``coupled'', understating true coupled cohabitation particularly among younger people. I smooth the married share across nearby ages exactly as in Section~\ref{sec:dirichlet} (the same age-local, effective-sample-size-weighted decay, but for a single Beta-distributed share rather than a Dirichlet-distributed composition vector), then winsorized to $[0.05,
0.95]$ so no cell relies entirely on one spatial signal:
\[
\hat\mu_{s,g,\ell,a} = \text{winsorize}_{[0.05,\,0.95]}\big(\text{smoothed
married share of cell } (s,g,\ell,a)\big).
\]

\subsection{PUMA Calibration Factors}

For each PUMA $p$ and cell $(s,g,\ell,a)$, define the uncalibrated LGBTQ
share of the state total projected in Section~\ref{sec:acs}, now computed
\emph{within} the cell rather than pooled across all cells:
\[
S_p(s,g,\ell,a) = \frac{\sum_{b \in p} \hat{E}[\text{count}_{s,g,\ell,a,b}]}
{\sum_{p'} \sum_{b \in p'} \hat{E}[\text{count}_{s,g,\ell,a,b}]}.
\]
The blended spatial share for cell $(s,g,\ell,a)$ is a weighted average of
the two signals above, using that cell's own coupled share as the weight:
\[
\rho_p(s,g,\ell,a) = \hat\mu_{s,g,\ell,a}\, \rho_p^{\text{couple}}
\;+\; (1-\hat\mu_{s,g,\ell,a})\, \rho_p^{\text{single}}(s,a).
\]
From here the tilt machinery is unchanged from the original single-signal
version, just re-indexed by cell. The blended signal is logit-transformed
and mean-centered relative to the cell's own uncalibrated distribution,
\[
\tilde{\sigma}_p(s,g,\ell,a)
=
\text{logit}(\rho_p(s,g,\ell,a))
\;-\;
\sum_{p'} S_{p'}(s,g,\ell,a)\,\text{logit}(\rho_{p'}(s,g,\ell,a)),
\]
and the tilt factor, calibrated target share, and calibration factor follow
exactly as before:
\begin{gather*}
T_p(s,g,\ell,a) = \exp\!\big(\gamma\, \tilde{\sigma}_p(s,g,\ell,a)\big),
\\[4pt]
S_p^*(s,g,\ell,a) = \frac{S_p(s,g,\ell,a)\, T_p(s,g,\ell,a)}
{\sum_{p'} S_{p'}(s,g,\ell,a)\, T_{p'}(s,g,\ell,a)},
\\[4pt]
\hat{c}_p(s,g,\ell,a) = \frac{S_p^*(s,g,\ell,a)}{S_p(s,g,\ell,a)}.
\end{gather*}
$\gamma$ keeps its original role as a single global conservative dial (site
still defaults to $\gamma=0.5$): it governs how strongly \emph{any} cell's
reweighting responds to its blended spatial signal, while the coupled-share
weight $\hat\mu$ is what determines \emph{which} spatial signal -- couples
or singles -- that cell is actually being tilted toward.

\subsection{Inverse-Distance-Weighted Smoothing to Tract Level}

Applying $\hat c_p(s,g,\ell,a)$ uniformly to every tract within PUMA $p$
produces a step-function surface, with hard discontinuities at PUMA
boundaries. Yet those jumps stem from where the Census Bureau happened to
draw PUMA lines, not real demographic transitions. Hence, I smooth the
calibration factors to the tract level using inverse-distance weighting
(IDW) instead of applying them as a step function.

Let $\mathbf{x}_b$ and $\mathbf{x}_p$ denote the centroids of tract $b$ and
PUMA $p$. The IDW-smoothed calibration factor for tract $b$ and cell
$(s,g,\ell,a)$ is
\[
\tilde{c}_b(s,g,\ell,a)
=
\frac{\displaystyle\sum_p \|\mathbf{x}_b - \mathbf{x}_p\|^{-2}\,
\hat{c}_p(s,g,\ell,a)}
     {\displaystyle\sum_p \|\mathbf{x}_b - \mathbf{x}_p\|^{-2}}.
\]
A tract deep inside a single PUMA receives a factor close to that PUMA's $\hat{c}_p$; a tract near a PUMA boundary receives a blend of the neighboring PUMAs' factors, weighted by inverse squared distance. IDW smoothing preserves geographic heterogeneity while producing a spatially continuous surface with no visible seams.

All subcategory estimates in tract $b$ and cell $(s,g,\ell,a)$ are
multiplied by $\tilde c_b(s,g,\ell,a)$ and then summed over $s,a$ and
collapsed over $g$ into the output bucket $k$, exactly as in
Section~\ref{sec:acs}; the whole state is then rescaled so totals match the
uncalibrated projection:
\[
\text{scale} =
\frac{\sum_b \hat{E}[\text{LGBTQ}_b]}{\sum_b \tilde c_b\, \hat{E}[\text{LGBTQ}_b]}.
\]

\section{Capacity Constraint}
\label{sec:capacity}

The PUMA reweighting above is a multiplicative tilt, which lacks a natural ceiling. In high-concentration areas, reweighting and state-level rescaling can push a tract's estimated count for an output gender bucket $k$ above that bucket's actual population. Yet no tract can contain more LGBTQ men than it contains men (or women or envies).

For output gender bucket $k \in \{m, w, \text{nb}\}$ and tract $b$, let
$\hat N^{\text{LGBTQ}}_{k,b}$ denote the post-calibration LGBTQ count and
$N_{k,b}$ the bucket's total population. I cap each bucket at its
population ceiling and scale every LGBTQ subcategory within that bucket down
proportionally, preserving relative composition:
\[
\kappa_{k,b} = \min\!\left(1,\; \frac{N_{k,b}}{\hat N^{\text{LGBTQ}}_{k,b}}
\right),
\]
with every subcategory count for bucket $k$ in tract $b$ multiplied by
$\kappa_{k,b}$. Imposition of the capacity constraint is rare. 

\end{document}
